


%----------------------------------------------------------------------------------------

\documentclass[oneside,openany]{tufte-book} % Use the tufte-book class which in turn uses the tufte-common class



\hypersetup{colorlinks} % Comment this line if you don't wish to have colored links

\usepackage{amsmath}
\usepackage{media9}
\usepackage{microtype} % Improves character and word spacing

\usepackage{lipsum} % Inserts dummy text

\usepackage{booktabs} % Better horizontal rules in tables
\hypersetup{pdfstartview=FitH, pdfpagemode=UseNone}
\usepackage{amsfonts,amssymb,graphics,epsfig,verbatim,bm,latexsym,amsmath,url,amsbsy}
\usepackage{graphicx} % Needed to insert images into the document
\graphicspath{{graphics/}} % Sets the default location of pictures
\setkeys{Gin}{width=\linewidth,totalheight=\textheight,keepaspectratio} % Improves figure scaling
\definecolor{mycolor}{rgb}{0.122, 0.435, 0.698}% Rule colour
\makeatletter
\newcommand{\mybox}[1]{%
	\setbox0=\hbox{#1}%
	\setlength{\@tempdima}{\dimexpr\wd0+13pt}%
	\begin{tcolorbox}[colframe=mycolor,boxrule=0.5pt,arc=4pt,
		left=6pt,right=6pt,top=6pt,bottom=6pt,boxsep=0pt,width=\@tempdima]
		#1
	\end{tcolorbox}
}
\makeatother
\usepackage{fancyvrb} % Allows customization of verbatim environments
\fvset{fontsize=\normalsize} % The font size of all verbatim text can be changed here

\definecolor{bananamania}{rgb}{0.98, 0.91, 0.71}
\definecolor{bananayellow}{rgb}{1.0, 0.88, 0.21}
\newcommand{\hangp}[1]{\makebox[0pt][r]{(}#1\makebox[0pt][l]{)}} % New command to create parentheses around text in tables which take up no horizontal space - this improves column spacing
\newcommand{\hangstar}{\makebox[0pt][l]{*}} % New command to create asterisks in tables which take up no horizontal space - this improves column spacing

\usepackage{xspace} % Used for printing a trailing space better than using a tilde (~) using the \xspace command

\newcommand{\highlight}[1]{\colorbox{green}{\ensuremath{#1}}}

\newcommand{\monthyear}{\ifcase\month\or January\or February\or March\or April\or May\or June\or July\or August\or September\or October\or November\or December\fi\space\number\year} % A command to print the current month and year

\newcommand{\openepigraph}[2]{ % This block sets up a command for printing an epigraph with 2 arguments - the quote and the author
	\begin{fullwidth}
		\sffamily\large
		\begin{doublespace}
			\noindent\allcaps{#1}\\ % The quote
			\noindent\allcaps{#2} % The author
		\end{doublespace}
	\end{fullwidth}
}
\geometry{
	left=22.8mm, % left margin
	textwidth=125mm, % main text block
	marginparsep=2.2mm, % gutter between main text block and margin notes
	marginparwidth=50.4mm % width of margin notes
}
\newcommand{\blankpage}{\newpage\hbox{}\thispagestyle{empty}\newpage} % Command to insert a blank page
\usepackage{hyperref}
\usepackage{tcolorbox}
\definecolor{indianyellow}{rgb}{0.89, 0.66, 0.34}
\definecolor{cadmiumyellow}{rgb}{1.0, 0.96, 0.0}
\tcbuselibrary{theorems}
\definecolor{beaublue}{rgb}{0.74, 0.83, 0.9}
\definecolor{aliceblue}{rgb}{0.94, 0.97, 1.0}
\definecolor{lightgoldenrodyellow}{rgb}{0.98, 0.98, 0.82}
\newtcbtheorem[number within=chapter]{myex}{Example}%
{colback=lightgoldenrodyellow,colframe=indianyellow,fonttitle=\bfseries}{th}
\newtcbtheorem[number within=chapter]{mydef}{Definition}%
{colback=beaublue,colframe=titlepagecolor,fonttitle=\bfseries}{th}





\usepackage{units} % Used for printing standard units

\newcommand{\hlred}[1]{\textcolor{Maroon}{#1}} % Print text in maroon
\newcommand{\hangleft}[1]{\makebox[0pt][r]{#1}} % Used for printing commands in the index, moves the slash left so the command name aligns with the rest of the text in the index
\newcommand{\hairsp}{\hspace{1pt}} % Command to print a very short space
\newcommand{\ie}{\textit{i.\hairsp{}e.}\xspace} % Command to print i.e.
\newcommand{\eg}{\textit{e.\hairsp{}g.}\xspace} % Command to print e.g.
\newcommand{\na}{\quad--} % Used in tables for N/A cells
\newcommand{\measure}[3]{#1/#2$\times$\unit[#3]{pc}} % Typesets the font size, leading, and measure in the form of: 10/12x26 pc.
\newcommand{\tuftebs}{\symbol{'134}} % Command to print a backslash in tt type in OT1/T1

\providecommand{\XeLaTeX}{X\lower.5ex\hbox{\kern-0.15em\reflectbox{E}}\kern-0.1em\LaTeX}
\newcommand{\tXeLaTeX}{\XeLaTeX\index{XeLaTeX@\protect\XeLaTeX}} % Command to print the XeLaTeX logo while simultaneously adding the position to the index

\newcommand{\doccmdnoindex}[2][]{\texttt{\tuftebs#2}} % Command to print a command in texttt with a backslash of tt type without inserting the command into the index

\newcommand{\doccmddef}[2][]{\hlred{\texttt{\tuftebs#2}}\label{cmd:#2}\ifthenelse{\isempty{#1}} % Command to define a command in red and add it to the index
	{ % If no package is specified, add the command to the index
		\index{#2 command@\protect\hangleft{\texttt{\tuftebs}}\texttt{#2}}% Command name
	}
	{ % If a package is also specified as a second argument, add the command and package to the index
		\index{#2 command@\protect\hangleft{\texttt{\tuftebs}}\texttt{#2} (\texttt{#1} package)}% Command name
		\index{#1 package@\texttt{#1} package}\index{packages!#1@\texttt{#1}}% Package name
	}}
	
	\newcommand{\doccmd}[2][]{% Command to define a command and add it to the index
		\texttt{\tuftebs#2}%
		\ifthenelse{\isempty{#1}}% If no package is specified, add the command to the index
		{%
			\index{#2 command@\protect\hangleft{\texttt{\tuftebs}}\texttt{#2}}% Command name
		}
		{%
			\index{#2 command@\protect\hangleft{\texttt{\tuftebs}}\texttt{#2} (\texttt{#1} package)}% Command name
			\index{#1 package@\texttt{#1} package}\index{packages!#1@\texttt{#1}}% Package name
		}}
		
		% A bunch of new commands to print commands, arguments, environments, classes, etc within the text using the correct formatting
		\newcommand{\docopt}[1]{\ensuremath{\langle}\textrm{\textit{#1}}\ensuremath{\rangle}}
		\newcommand{\docarg}[1]{\textrm{\textit{#1}}}
		\newenvironment{docspec}{\begin{quotation}\ttfamily\parskip0pt\parindent0pt\ignorespaces}{\end{quotation}}
		\newcommand{\docenv}[1]{\texttt{#1}\index{#1 environment@\texttt{#1} environment}\index{environments!#1@\texttt{#1}}}
		\newcommand{\docenvdef}[1]{\hlred{\texttt{#1}}\label{env:#1}\index{#1 environment@\texttt{#1} environment}\index{environments!#1@\texttt{#1}}}
		\newcommand{\docpkg}[1]{\texttt{#1}\index{#1 package@\texttt{#1} package}\index{packages!#1@\texttt{#1}}}
		\newcommand{\doccls}[1]{\texttt{#1}}
		\newcommand{\docclsopt}[1]{\texttt{#1}\index{#1 class option@\texttt{#1} class option}\index{class options!#1@\texttt{#1}}}
		\newcommand{\docclsoptdef}[1]{\hlred{\texttt{#1}}\label{clsopt:#1}\index{#1 class option@\texttt{#1} class option}\index{class options!#1@\texttt{#1}}}
		\newcommand{\docmsg}[2]{\bigskip\begin{fullwidth}\noindent\ttfamily#1\end{fullwidth}\medskip\par\noindent#2}
		\newcommand{\docfilehook}[2]{\texttt{#1}\index{file hooks!#2}\index{#1@\texttt{#1}}}
		\newcommand{\doccounter}[1]{\texttt{#1}\index{#1 counter@\texttt{#1} counter}}
		
		
		
		\titleformat{\chapter}%
		[block]% shape
		{\relax\ifthenelse{\NOT\boolean{@tufte@symmetric}}{\begin{fullwidth}}{}}% format applied to label+text
			{\itshape\huge\thechapter\hspace{10pt}}% label
			{0pt}% horizontal separation between label and title body
			{\huge\rmfamily\itshape}% before the title body
			[\ifthenelse{\NOT\boolean{@tufte@symmetric}}{\end{fullwidth}}{}]% after the title body
		
		
		
		
		\hypersetup{
			colorlinks,
			citecolor=green,
			filecolor=black,
			linkcolor=blue,
			urlcolor=red
		}
		
		\setcounter{secnumdepth}{3}
		
		\setcounter{tocdepth}{2}
		
		
		\definecolor{titlepagecolor}{cmyk}{1,.60,0,.40}
		\definecolor{ikb}{rgb}{0.0, 0.18, 0.65}
		\definecolor{namecolor}{cmyk}{1,.50,0,.10}
		
		\usepackage{makeidx} % Used to generate the index
		\makeindex % Generate the index which is printed at the end of the document
		
		
		\begin{document}
			\title{ECON5221: Econometric Time Series Analysis }
	
			
			\begin{titlepage}
				\newgeometry{left=7.5cm} %defines the geometry for the titlepage
				\pagecolor{ikb}
				\noindent
				
				\color{white}
				\makebox[0pt][l]{\rule{1\textwidth}{2pt}}
				\par
				
				\vspace{5pt}\textbf{\huge \textsf{Time Series Econometrics}} \hspace{5pt} \textcolor{white}{\textsf{Nicky Grant}}
				\vfill
				\noindent
				{\huge \textsf{ECON5221: Problem Set 1}}
				\vskip\baselineskip
				\noindent
				\textsf{\Large Semester 2 2020/2021}
			\end{titlepage}
			\restoregeometry % restores the geometry
			\nopagecolor% Use this to restore the color pages to whiteint the title page
			
			







\begin{fullwidth}





\section*{\LARGE Exercise Questions 1}


\emph{The questions in Problem Set 1  \ref{c1} are to be covered in Tutorial 1 (as many as we can get through in an hour). The questions in Video Problem Set 1 are to be completed by you also,  a full video solution will be given after the first tutorial.}\newline

\vspace{20pt}
\subsection*{Problem Set 1}\label{c1}
\begin{enumerate}
	\item[] \textbf{In the following questions each process evolves over time }$\mathcal{T}=\{\cdots, -2, -1,0,1,2,\cdots\}$.

\vspace{15pt}
	
	\item An $MA(2)$ process is%
	\[
	Y_{t}=\alpha +\varepsilon _{t}+\theta _{1}\varepsilon _{t-1}+\theta
	_{2}\varepsilon _{t-2} \qquad 	 \varepsilon _{t}\sim WN(\sigma^2)
	\]

	\vspace{1cm}
	
	\begin{enumerate}
		\item Show that
		\[
		\mathbb{E}[Y_{t}]=\alpha
		\]%
		and%
		\[
		Var[Y_{t}]=\sigma ^{2}(1+\theta _{1}^{2}+\theta _{2}^{2}).
		\]
		
		\item  \noindent Find the autocorrelation function of $Y_{t}$.
		These autocorrelations should be written in terms of the $MA$ coefficients ($%
		\theta _{1}$ and $\theta _{2}$).$\smallskip $
		
	
	\end{enumerate}
	

	
	
	\vspace{18pt}
	
	
	\item Consider the AR(2) process
	\begin{equation} Y_t=\mu+ \phi_1 Y_{t-1}+\phi_2Y_{t-2} +\varepsilon_t  \qquad 	 \varepsilon _{t}\sim WN(\sigma^2) \end{equation}
	
	\begin{enumerate}
		\item Under what conditions on $\phi_1,\phi_2$ is the AR(2) process in (1) weakly stationary.
				\vspace{6pt}
		\item  Assuming conditions in 2(a) hold, derive the mean of $Y_t$.
				\vspace{6pt}
		\item Assuming conditions in 2(a) holds, derive a second order difference equation for the autocovariance function of $Y_t$.
		
		\item  Using 2(c) derive the variance of $Y_t$ and the correlation of $Y_t$ and $Y_{t-1}$.
	\end{enumerate}
	
	
	\vspace{18pt}
	 
	\item Consider an i.i.d process $\{Y_t\}$ where $Y_t$ is Cauchy distributed with  density function $f_{Y_t}(y)= \frac{1}{\pi(1+y^2)}$ for y on the real line $\mathbb{R}$. Show that the process is strictly stationary. Is the process also weakly stationary? [Provide a formal justification for your answer.] 
	
	\vspace{18pt}
	\item If a process $\{Y_T\}$ is weakly stationary with absolutely summable covariances which conditions on its first and second moments does it satisfy? Prove that $\bar{Y}_T=\frac{1}{T} \sum_{t=1}^TY_t$ is a consistent estimator of the mean of $Y_t$ under these conditions. [\textit{Hint: Use Chebychev's inequality}.]
	


\end{enumerate}
\subsection*{Video Problem Set 1}

\begin{enumerate}
	
	\item Consider an $MA(\infty )$ process%
	\[
	Y_{t}=\mu +\sum\limits_{s=0}^{\infty }\eta _{s}\varepsilon _{t-s}
	\]%
	where the $MA$ coefficients satisfy the difference equation%
	\[
	\eta _{s}=\phi _{1}\eta _{s-1},\quad s>1
	\]%
	with starting values
	\[
	\eta _{0}=1,\quad \eta _{1}=\phi _{1}+\theta _{1}.
	\]%

	
	\begin{enumerate}
		\item
		Obtain $Y_{t}-\phi _{1}Y_{t-1}$ and hence show that the process with this $%
		MA(\infty )$ representation is the $ARMA(1,1)$%
		\[
		Y_{t}=\alpha +\phi _{1}Y_{t-1}+\varepsilon _{t}+\theta _{1}\varepsilon _{t-1}
		\]%
		with $\alpha =(1-\phi _{1})\mu $.\medskip
		
		\item Using the difference equation for the $MA(\infty )$ coefficients, show
		that%
		\[
		\eta _{s}=\phi _{1}^{s-1}\eta _{1},\quad s>1
		\]%
		with $\eta _{1}=\phi _{1}+\theta _{1}.\medskip $
		
		\item The absolute summablity condition states that this process is
		stationary provided%
		\[
		\sum\limits_{s=0}^{\infty }\left\vert \eta _{s}\right\vert <\infty .
		\]
		
		Using the result from (b), show that the $ARMA(1,1)$ process is stationary
		when%
		\[
		\left\vert \phi _{1}\right\vert <1.
		\]%
		[Hint: Note that for any two values $a$ and $b$, $\left\vert ab\right\vert
		=\left\vert a\right\vert \times \left\vert b\right\vert $.]$\medskip $
		
		\item Again using the result from (b), together with the general result that
		the variance of the $MA(\infty )$ process is given by%
		\[
		Var[Y_{t}]=\sigma ^{2}\sum\limits_{s=0}^{\infty }\eta _{s}^{2},
		\]%
		show that the variance of the $ARMA(1,1)$ process is%
		\begin{eqnarray*}
			Var[Y_{t}] &=&\sigma ^{2}\left\{ 1+(\phi _{1}+\theta
			_{1})^{2}\sum\limits_{s=0}^{\infty }\phi _{1}^{2s}\right\}  \\
			&=&\sigma ^{2}\frac{(1+2\phi _{1}\theta _{1}+\theta _{1}^{2})}{1-\phi
				_{1}^{2}}.
		\end{eqnarray*}
		
		\item Show that \[\gamma(k) =\left\{
		\begin{array}{lr}
		\sigma^{2}(\phi_1+\theta_1)\left[1+\frac{(\phi_1+\theta_1)\phi_1}{1-\phi_1^2} \right]   & : k =1\\
		\phi_1^{k-1}\gamma(1) & : k>1
		\end{array}
		\right.
		\]
		
		
		
		\[\rho(k) =\left\{
		\begin{array}{lr}
		\frac{(\phi_1+\theta_1)(1+\phi_1\theta_1)}{1+2\phi_1\theta_1+\theta_1^2}   & : k =1\\
		\phi_1^{k-1}\rho(1) & : k>1
		\end{array}
		\right.
		\]
		
		\item What happens to the $MA(\infty )$ coefficients $\eta _{s}$ ($s=1,2,...)
		$ given above in the special case where $\theta _{1}=-\phi _{1}$.  Comment on the result.
	\end{enumerate}
\end{enumerate}
\end{fullwidth}
\end{document}